% LaTeX master Datei(en) zusammengestellt von Falk-Jonatan Strube zur Nutzung an der Hochschule für Technik und Wirtschaft Dresden: https://github.com/genericFJS/htw
\documentclass{scrreprt}
\gdef\pathtomaster{../_LaTeX_master}
%% HTW Dresden Corporate Design
%% Copyright (C) Falk-Jonatan Strube <falkjonatanstrube@gmail.com>, 2017
%%% === Anfang Paket htwscr ===
\makeatletter
\ifcsname KOMAoption\endcsname%
\KOMAoption{footheight}{11mm}	% mehr Platz für Wasserzeichen o.ä. in der Fußzeile
\recalctypearea	% s.o.
\RequirePackage{geometry}
\geometry{top=25mm, bottom=25mm, headsep=10mm, inner=25mm, outer=25mm, footskip=15mm, marginparsep=2mm}
\RequirePackage{helvet}	% Corporate Design schreibt Arial/Helvetica vor
\renewcommand{\familydefault}{\sfdefault}	% s.o.
\RequirePackage{graphicx}	% für Grafiken
\RequirePackage{xcolor}		% für Farben
\definecolor{htworange}{RGB}{249,155,28}	% Corporate Design orange CMYK(0/45/100/0)
\definecolor{htwblue}{RGB}{0,116,188}			% Corporate Design blau CMYK(90/50/0/0)
\definecolor{htwgrey}{RGB}{128,128,128}		% Corporate Design gray 60% Schwarz [Annahme: CMYK(0,0,0,60)]
\RequirePackage{ifthen,xifthen}	% für Bedingungen/Optionen
\RequirePackage{etoolbox}				% s.o.

%%% Format des HTW-Headers (wird hier als Texteelement verwenden) nach dem Style der PowerPoint-Vorlagen. Weiterhin inspiriert durch das Paket tudscr (Corporate Design der TU Dresden)
%%% Dateipfad des Logos abhängig von der Verwendung als Paket/Input
\ifdefined\@isHTWCDpackage%
	\gdef\@pathtoimage{HTW}%
\else%
	\ifdefined\pathtomaster%
		\gdef\@pathtoimage{\pathtomaster /HTW}%
	\else 
		\gdef\@pathtoimage{HTW}%
	\fi%
\fi%
\newcommand{\pathtologo}[1]{\gdef\@pathtoimage{#1}}
%%% Berechnung der Größe des Logos
\newlength{\htwLogoHeight}
\setlength{\htwLogoHeight}{12mm}	% Mindesthöhe des Logos laut Corporate Design: 11mm
\newlength{\@htwLogoWidth}
\AtBeginDocument{
\def\printLogo{\includegraphics[height=\htwLogoHeight]{\@pathtoimage}}
\setlength\@htwLogoWidth{\widthof{\printLogo}}
}
%%% Definition des Header
\newcommand{\@htwheader}{{\noindent%
%%% Header ist abhängig von dem Vorhandensein von department, institute und chair
\newcommand{\hasDepInsCh}{\gdef\@hasDepInsCh{}}%
\ifdefined\@department\hasDepInsCh\fi%
\ifdefined\@institute\hasDepInsCh\fi%
\ifdefined\@chair\hasDepInsCh\fi%
\begin{minipage}[b][\htwLogoHeight]{\dimexpr \textwidth - \@htwLogoWidth - 2mm \relax}% 2mm Abstand zwischen Logo und Linie
{\small{\bfseries%
\ifdefined\@faculty\@faculty{}\fi} %
\strut}%
{\color{htworange}\hrule height 1pt}%
\vspace{.1mm}%
\end{minipage}
\hfill%
\begin{minipage}[b][\htwLogoHeight]{\@htwLogoWidth}
\includegraphics[height=\htwLogoHeight]{\@pathtoimage}
\end{minipage}
\ifdefined\@hasDepInsCh%
\\%
\begin{minipage}{\textwidth}
\small%
\ifdefined\@department\@department{}\fi%
\ifdefined\@institute\ifdefined\@department,~\fi \@institute\fi%
\ifdefined\@chair\ifdefined\@department,~\else\ifdefined\@institute,~\fi\fi\@chair\fi\strut
%{\color{htworange}\hrule height 0.5pt}
\end{minipage}
\fi
}}

%%% Definition der Fußzeilen
\usepackage{scrlayer-scrpage}
\@ifclassloaded{scrreprt}{%
\renewcommand*{\partpagestyle}{scrheadings}
\renewcommand*{\chapterpagestyle}{scrheadings}%
}{}
\pagestyle{scrheadings}
\cfoot{\vfill\pagemark}
\chead{}

%%% Einstellungen
\newcommand{\simpledate}{\gdef\@simpledate{}}		% Datum nicht mit "Eingereicht am:" Vorsatz
\newcommand{\nouppercase}{\gdef\@nouppercase{}}	% Titel/Parts/Chapters usw. nicht in Großbuchstaben

%%% Einiges erst zu Beginn des Dokumentes ausführen, um Einstellungen zu berücksichtigen
\AtBeginDocument{
\renewcommand\SS{SS}	% fix big ß
%%% Sektionen usw. ggf. in Großbuchstaben:
\providecommand{\chapterlinesformat}[3]{}	% compatibility
\providecommand{\chapterlineswithprefixformat}[3]{}	% compatibility
\providecommand{\sectionlinesformat}[4]{}	% compatibility
\ifdefined\@nouppercase\relax\else
\renewcommand\sectionlinesformat[4]{\@hangfrom{\hskip#2 #3}{\MakeUppercase{#4}}}
\@ifclassloaded{scrreprt}{%
\renewcommand\chapterlinesformat[3]{\@hangfrom{#2}{\MakeUppercase{#3}}}
\renewcommand\chapterlineswithprefixformat[3]{\MakeUppercase{#2#3}}%
}{}
\fi
%%% Anpassung der Part-Seite
\renewcommand*{\partheadstartvskip}{\@htwheader\vspace*{30mm}}
\renewcommand*{\raggedpart}{\raggedright}
\renewcommand*{\partformat}{\huge\ifdefined\@nouppercase \partname~\thepart\else\MakeUppercase{\partname~\thepart}\fi.}
\ifdefined\@nouppercase\setkomafont{part}{\Huge}\else\setkomafont{part}{\Huge\MakeUppercase}\fi
}

\newcommand{\bilingual}[3]{\ifdefined #1 \let #1\undefined\fi\newcommand{#1}{\iflanguage{english}{#3}{#2}}}	% Übersetzung von einigen Variablen in mehrere Sprachen

\newcommand{\checkThesis}[2]{\ifthenelse{\equal{#1}{diss}\OR\equal{#1}{doctoral}\OR\equal{#1}{phd}}{
\gdef#2{\dissertationname}}{
\ifthenelse{\equal{#1}{diploma}}{
\gdef#2{\diplomathesisname}}{
\ifthenelse{\equal{#1}{master}}{
\gdef#2{\masterthesisname}}{
\ifthenelse{\equal{#1}{bachelor}}{
\gdef#2{\bachelorthesisname}}{
\ifthenelse{\equal{#1}{student}}{
\gdef#2{\studentthesisname}}{
\ifthenelse{\equal{#1}{evidence}}{
\gdef#2{\studentresearchname}}{
\ifthenelse{\equal{#1}{project}}{
\gdef#2{\projectpapername}}{
\ifthenelse{\equal{#1}{seminar}}{
\gdef#2{\seminarpapername}}{
\ifthenelse{\equal{#1}{term}}{
\gdef#2{\termpapername}}{
\ifthenelse{\equal{#1}{research}}{
\gdef#2{\researchname}}{
\ifthenelse{\equal{#1}{log}}{
\gdef#2{\logname}}{
\ifthenelse{\equal{#1}{report}}{
\gdef#2{\reportname}}{
\ifthenelse{\equal{#1}{internship}}{
\gdef#2{\internshipname}}{
\ifthenelse{\equal{#1}{lecture}\OR\equal{#1}{lesson}\OR\equal{#1}{pract}}{
\bilingual{\professorname}{Vorlesung von}{Lecture by}
\simpledate
\gdef\@islecture{\relax}
\ifthenelse{\equal{#1}{lecture}}{\gdef#2{\lecturename}}{
\ifthenelse{\equal{#1}{lesson}}{\gdef#2{\lessonname}}{
\gdef#2{\practname}
}}}{
\gdef#2{#1}}}}}}}}}}}}}}}}

%%% Textfragmente der Titelseite in deutsch und englisch
\bilingual{\dateofbirthname}{Geboren am:}{Born on:}
\bilingual{\placeofbirthname}{in}{in}
\bilingual{\coursename}{Studiengang:}{Course:}
\bilingual{\disciplinename}{Studienrichtung:}{Discipline:}
\bilingual{\matriculationnumbername}{Matrikelnummer:}{Matriculation number:}
\bilingual{\matriculationyearname}{Immatrikulationsjahr}{Matriculation year:}
\bilingual{\graduationname}{zur Erlangung des akademischen Grades}{to achieve the academic degree}
\bilingual{\refereename}{Gutachter}{Referee}
\bilingual{\advisorname}{Fachreferent}{Advisor}
\bilingual{\supervisorname}{Betreuer}{Supervisor}
\bilingual{\professorname}{Betreuender Hochschullehrer}{Supervising professor}
\bilingual{\datename}{Eingereicht am:}{Submitted on:}
%% Thesis Typen (inspiriert von dem Paket tudscr):
\bilingual{\dissertationname}{Dissertation}{Dissertation} % diss, doctoral, phd 
\bilingual{\diplomathesisname}{Diplomarbeit}{Diploma Thesis} % diploma 
\bilingual{\masterthesisname}{Master-Arbeit}{Master Thesis} % master 
\bilingual{\bachelorthesisname}{Bachelor-Arbeit}{Bachelor Thesis} % bachelor 
\bilingual{\studentthesisname}{Studienarbeit}{Student Thesis} % student 
\bilingual{\studentresearchname}{Großer Beleg}{Student Research Project} % evidence 
\bilingual{\projectpapername}{Projektarbeit}{Project Paper} % project
\bilingual{\seminarpapername}{Seminararbeit}{Seminar Paper} % seminar
\bilingual{\termpapername}{Hausarbeit}{Term Paper} % term
\bilingual{\researchname}{Forschungsbericht}{Research Report} % research 
\bilingual{\logname}{Protokoll}{Log} % log 
\bilingual{\reportname}{Bericht}{Report} % report
\bilingual{\internshipname}{Praktikumsbericht}{Internship Report} % internship 
\bilingual{\lecturename}{Vorlesungsmitschrift}{Lecture Notes} % lecture 
\bilingual{\lectureauthorname}{Mitschrift von}{Notes by} % lecture  
\bilingual{\lessonname}{\"Ubungsmitschrift}{Lesson Notes} % lecture  
\bilingual{\practname}{Praktikumsmitschrift}{Pracital Course Notes} % lecture 
%%% Variablen für die Titelseite (inspiriert von dem Paket tudscr):
\renewcommand{\subject}[1]{\checkThesis{#1}{\@subject}}
\newcommand{\faculty}[1]{\gdef\@faculty{#1}}
\newcommand{\department}[1]{\gdef\@department{#1}}
\newcommand{\institute}[1]{\gdef\@institute{#1}}
\newcommand{\chair}[1]{\gdef\@chair{#1}}
\newcommand{\authormore}[1]{\gdef\@authormore{#1}}
\newcommand{\dateofbirth}[1]{\gdef\@dateofbirth{#1}}
\newcommand{\placeofbirth}[1]{\gdef\@placeofbirth{#1}}
\newcommand{\course}[1]{\gdef\@course{#1}}
\newcommand{\discipline}[1]{\gdef\@discipline{#1}}
\newcommand{\matriculationnumber}[1]{\gdef\@matriculationnumber{#1}}
\newcommand{\matriculationyear}[1]{\gdef\@matriculationyear{#1}}
\newcommand{\thesis}[1]{\checkThesis{#1}{\@thesis}}
\newcommand{\graduation}[2][]{\gdef\@graduation{#2}\gdef\@graduationshort{#1}}
\newcommand{\supervisor}[1]{\gdef\@supervisor{#1}}
\newcommand{\referee}[1]{\gdef\@referee{#1}}
\newcommand{\advisor}[1]{\gdef\@advisor{#1}}
\newcommand{\professor}[1]{\gdef\@professor{#1}}
%%% Titelseite (inspiriert von dem Paket tudscr):
\renewcommand*{\maketitle}{%
\begin{titlepage}%
\begin{flushleft}
\renewcommand{\thefootnote}{$\star$}%
%%% Kopfzeile HTW-Style:
\@htwheader\\[13mm]
%%% Die Zeilen für titlehead, subject, title, subtitle und author sind immer vorhanden (spacing)
{\noindent\@titlehead{}\strut}\\[3mm]%
{\Large\@subject{}\strut}\\[8mm]%
{\Huge\bfseries%
\ifdefined\@nouppercase\@title{}\else\MakeUppercase{\@title{}}\fi%
\strut}\\
{\Large\bfseries%
\ifdefined\@nouppercase\@subtitle{}\else\MakeUppercase{\@subtitle{}}\fi%
\strut}\\[11mm]
\ifdefined\@islecture\strut\else{\Large\@author{}\strut}\fi % Autor wird bei Vorlesungsmitschrift unten angezeigt
%%% Weitere Felder werden nur angezeigt, wenn sie definiert wurden
\ifdefined\@authormore\ifdefined\@islecture\else\\\fi{\@authormore{}\strut}\fi
\ifdefined\@dateofbirth\\\dateofbirthname{} \@dateofbirth{} \ifdefined\@placeofbirth\placeofbirthname{} \@placeofbirth{}\strut\fi\fi
\ifdefined\@course\\\coursename{} \@course{}\strut\fi
\ifdefined\@discipline\\\disciplinename{} \@discipline{}\strut\fi
\ifdefined\@matriculationnumber\\\matriculationnumbername{} \@matriculationnumber{}\strut\fi
\ifdefined\@matriculationyear\\\matriculationyearname{} \@matriculationyear{}\strut\fi
\ifdefined\@thesis\\[10mm]{\LARGE\bfseries
\ifdefined\@nouppercase\@thesis{}\else\MakeUppercase{\@thesis{}}\fi
\strut}
\ifdefined\@graduation\\[2mm]{\graduationname{}\strut}
\\[5mm]{\LARGE\bfseries
\ifdefined\@nouppercase\relax\else\MakeUppercase{\fi%
\@graduation{} \ifdefined\@graduationshort(\@graduationshort{})\fi
\ifdefined\@nouppercase\relax\else}\fi\strut}\fi\fi
\vfill\noindent
\ifdefined\@islecture\\[1mm]{\lectureauthorname{}\strut}\\[.7mm]{\Large\@author{}\strut}\fi
\ifdefined\@referee{\refereename{}\strut}\\[.7mm]{\Large\@referee{}\strut}\fi
\ifdefined\@advisor\\[1mm]{\advisorname{}\strut}\\[.7mm]{\Large\@advisor{}\strut}\fi
\ifdefined\@supervisor\\[1mm]{\supervisorname{}\strut}\\[.7mm]{\Large\@supervisor{}\strut}\fi
\ifdefined\@professor\\[1mm]{\professorname{}\strut}\\[.7mm]{\Large\@professor{}\strut}\fi
\ifdefempty{\@date}{}{\ifdefined\@simpledate\\[6mm]\relax\else\\[3mm]\datename{}~\fi\@date{}\strut}
\ifdefempty{\@publishers}{}{\\[4.5mm]{\Large\@publishers{}\strut}}%
\renewcommand\footnoterule{}
\@thanks\global\let\@thanks\@empty
\end{flushleft}
\end{titlepage}
\setcounter{footnote}{0}
}
\else%
\PackageError{htwcd}{htwcd needs a KOMA class!}%
\fi%
\makeatother
%%% === Ende Paket htwscr ===
%-- Nummerierung bis Subsubsection für Report
\setcounter{secnumdepth}{3}	% zählt auch subsubsection
\setcounter{tocdepth}{3}	% Inhaltsverzeichnis bis in subsubsection

% LuaLaTeX Unterscheidung
\usepackage{ifluatex}

% Input inkl. Umlaute, Silbentrennung
\ifluatex
	\usepackage{fontspec}
\else
	\usepackage[utf8]{inputenc}
	\usepackage[T1]{fontenc}
\fi
\usepackage[english,ngerman]{babel}
\usepackage{csquotes}	% Anführungszeichen
\usepackage{marvosym}
\usepackage{eurosym}

% Style-Aufhübschung
\usepackage{soul, color}	% Kapitälchen, Unterstrichen, Durchgestrichen usw. im Text
%\usepackage{titleref}
\usepackage[breakwords, fit]{truncate}	% Abschneiden von Sätzen
\renewcommand{\TruncateMarker}{\,…}

% Mathe usw.
\usepackage{amssymb}
\usepackage{amsthm}
\usepackage[fleqn,intlimits]{amsmath}	% fleqn: align-Umgebung rechtsbündig; intlimits: Integralgrenzen immer ober-/unterhalb
%\usepackage{mathtools} % u.a. schönere underbraces
\usepackage{xcolor}
\usepackage{esint}	% Schönere Integrale, \oiint vorhanden
\everymath=\expandafter{\the\everymath\displaystyle}	% Mathe Inhalte werden weniger verkleinert
\usepackage{wasysym}	% mehr Symbole, bspw \lightning
%\renewcommand{\int}{\int\limits}
%\usepackage{xfrac}	% mehr fracs: sfrac{}{}
\let\oldemptyset\emptyset	% schöneres emptyset
\let\emptyset\varnothing
%\RequirePackage{mathabx}	% mehr Symbole
\mathchardef\mhyphen="2D	% Hyphen in Math
\usepackage{mathtools}		% zum Beispiel {pmatrix}[r] für Formatierung der Matrix

% tikz usw.
\usepackage{tikz}
\usepackage{pgfplots}
\pgfplotsset{compat=1.11}	% Umgeht Fehlermeldung
\usetikzlibrary{graphs}
%\usetikzlibrary{through}	% ???
\usetikzlibrary{arrows}
\usetikzlibrary{arrows.meta}	% Pfeile verändern / vergrößern: \draw[-{>[scale=1.5]}] (-3,5) -> (-3,3);
\usetikzlibrary{automata,positioning} % Zeilenumbruch im Node node[align=center] {Text\\nächste Zeile} automata für Graphen
\usetikzlibrary{matrix}
\usetikzlibrary{patterns}	% Schraffierte Füllung
\usetikzlibrary{shapes.geometric}	% Polygon usw.
\tikzstyle{reverseclip}=[insert path={	% Inverser Clip \clip
	(current page.north east) --
	(current page.south east) --
	(current page.south west) --
	(current page.north west) --
	(current page.north east)}
% Nutzen: 
%\begin{tikzpicture}[remember picture]
%\begin{scope}
%\begin{pgfinterruptboundingbox}
%\draw [clip] DIE FLÄCHE, IN DER OBJEKT NICHT ERSCHEINEN SOLL [reverseclip];
%\end{pgfinterruptboundingbox}
%\draw DAS OBJEKT;
%\end{scope}
%\end{tikzpicture}
]	% Achtung: dafür muss doppelt kompliert werden!
\usepackage{graphpap}	% Grid für Graphen
\tikzset{every state/.style={inner sep=2pt, minimum size=2em}}
\usetikzlibrary{mindmap, backgrounds}
%\usepackage{tikz-uml}	% braucht Dateien: http://perso.ensta-paristech.fr/~kielbasi/tikzuml/

% Tabular
\usepackage{longtable}	% Große Tabellen über mehrere Seiten
\usepackage{multirow}	% Multirow/-column: \multirow{2[Anzahl der Zeilen]}{*[Format]}{Test[Inhalt]} oder \multicolumn{7[Anzahl der Reihen]}{|c|[Format]}{Test2[Inhalt]}
\renewcommand{\arraystretch}{1.3} % Tabellenlinien nicht zu dicht
\usepackage{colortbl}
\arrayrulecolor{gray}	% heller Tabellenlinien
\usepackage{array}	% für folgende 3 Zeilen (für Spalten fester breite mit entsprechender Ausrichtung):
\newcolumntype{L}[1]{>{\raggedright\let\newline\\\arraybackslash\hspace{0pt}}m{\dimexpr#1\columnwidth-2\tabcolsep-1.5\arrayrulewidth}}
\newcolumntype{C}[1]{>{\centering\let\newline\\\arraybackslash\hspace{0pt}}m{\dimexpr#1\columnwidth-2\tabcolsep-1.5\arrayrulewidth}}
\newcolumntype{R}[1]{>{\raggedleft\let\newline\\\arraybackslash\hspace{0pt}}m{\dimexpr#1\columnwidth-2\tabcolsep-1.5\arrayrulewidth}}
\usepackage{booktabs}	% Kompatibilität für array und \hrule
\usepackage{caption}	% Um auch unbeschriftete Captions mit \caption* zu machen

% Nützliches
\usepackage{verbatim}	% u.a. zum auskommentieren via \begin{comment} \end{comment}
\usepackage{tabto}	% Tabs: /tab zum nächsten Tab oder /tabto{.5 \CurrentLineWidth} zur Stelle in der Linie
\NumTabs{6}	% Anzahl von Tabs pro Zeile zum springen
\usepackage{listings} % Source-Code mit Tabs
\usepackage{lstautogobble} 
\usepackage{enumitem}	% Anpassung der enumerates
%\setlist[enumerate,1]{label=(\arabic*)}	% global andere Enum-Items
\renewcommand{\labelitemiii}{$\scriptscriptstyle ^\blacklozenge$} % global andere 3. Item-Aufzählungszeichen
\usepackage{letltxmacro} % neue Definiton von Grundbefehlen
% Nutzen:
%\LetLtxMacro{\oldemph}{\emph}
%\renewcommand{\emph}[1]{\oldemph{#1}}
\RequirePackage{xpatch}	% ua. Konkatenieren von Strings/Variablen (etoolbox)
\usepackage{xstring}	% String Operationen
\usepackage{minibox}	% Minibox anstatt \fbox{} für Boxen mit (manuellen) Zeilenumbruch
\usepackage{framed}		% \begin{framed} statt minibox/fbox mit automatischem Zeilenumbruch
\usepackage{microtype}	% Verfeinerte Hyphenation
\usepackage{import}		
\usepackage{standalone}	% Zusammen mit import: Einfügen von separaten Latex-Dateien (die schon header haben)

% Einrichtung von lst
%\renewcommand*{\ttdefault}{qcr}
\lstset{
basicstyle=\ttfamily, 
%basicstyle=\footnotesize\sffamily\color{black},
%mathescape=true, 
%escapeinside=^^, 
autogobble, 
tabsize=2,
frame=single,
rulecolor=\color{lightgray},
numbers=left,
numbersep=5pt,
numberstyle=\tiny\color{gray},
commentstyle=\color{gray},
keywordstyle=\color{htwblue},
stringstyle=\color{htworange},
morecomment=[l][\color{magenta}]{\#}
showspaces=false,
showstringspaces=false,
breaklines=true,
literate=%
    {Ö}{{\"O}}1
    {Ä}{{\"A}}1
    {Ü}{{\"U}}1
    {ß}{{\ss}}1
    {ü}{{\"u}}1
    {ä}{{\"a}}1
    {ö}{{\"o}}1
    {„}{{\glqq}}1
    {“}{{\grqq}}1
    {~}{{\textasciitilde}}1
}
\usepackage{scrhack} % Fehler umgehen
\def\ContinueLineNumber{\lstset{firstnumber=last}} % vor lstlisting. Zum wechsel zum nicht-kontinuierlichen muss wieder \StartLineAt1 eingegeben werden
\def\StartLineAt#1{\lstset{firstnumber=#1}} % vor lstlisting \StartLineAt30 eingeben, um bei Zeile 30 zu starten
\let\numberLineAt\StartLineAt

% BibTeX
\usepackage[
backend=bibtex,
bibencoding=ascii,
style=authoryear, citestyle=authoryear-ibid,
%doi=false,
%isbn=false,
%url=false
]{biblatex}	% BibTeX
\usepackage{makeidx}
%\makeglossary
%\makeindex
\usepackage{url}	% Darstellung von URLs

% Grafiken
\usepackage{graphicx}
\usepackage{epstopdf}	% eps-Vektorgrafiken einfügen
\usepackage{transparent}	% transparent nutzen: {\transparent{0.4} ...}
%\epstopdfsetup{outdir=\customDir}
% Prüft, ob Grafik existiert (mit \ifvalidimage{}{}) [Quelle: https://tex.stackexchange.com/a/99176]:
\makeatletter
\newif\ifgraphicexist
\catcode`\*=11
\newcommand\ifvalidimage[1]{%
    \begingroup
    \global\graphicexisttrue
    \let\input@path\Ginput@path
    \filename@parse{#1}%
    \ifx\filename@ext\relax
    \@for\Gin@temp:=\Gin@extensions\do{%
        \ifx\Gin@ext\relax
        \Gin@getbase\Gin@temp
        \fi}%
    \else
    \Gin@getbase{\Gin@sepdefault\filename@ext}%
    \ifx\Gin@ext\relax
    \global\graphicexistfalse
    \def\Gin@base{\filename@area\filename@base}%
    \edef\Gin@ext{\Gin@sepdefault\filename@ext}%
    \fi
    \fi
    \ifx\Gin@ext\relax
    \global\graphicexistfalse
    \else 
    \@ifundefined{Gin@rule@\Gin@ext}%
    {\global\graphicexistfalse}%
    {}%
    \fi  
    \ifx\Gin@ext\relax 
    \gdef\imageextension{unknown}%
    \else
    \xdef\imageextension{\Gin@ext}%
    \fi 
    \endgroup 
    \ifgraphicexist
    \expandafter \@firstoftwo
    \else
    \expandafter \@secondoftwo
    \fi 
} 
\catcode`\*=12
\makeatother
\LetLtxMacro{\forceincludegraphics}{\includegraphics}	% neuer Befehl für includegraphics
\renewcommand{\includegraphics}[2][]{	% altes includegraphics neu definieren, damit es auch nicht vorhandene einfügt
\ifvalidimage{#2}{
\forceincludegraphics[#1]{#2}
}{
\message{Achtung: Grafik wurde nicht gefunden: '#2'}
\minibox[frame]{
\textbf{\StrSubstitute{#2}{_}{\_}}%  \ifthenelse{\isempty{#1}}{}{\\\textit{#1}}
}}}

% pdf-Setup
\usepackage{pdfpages}
\usepackage[bookmarks,%
bookmarksopen=false,% Klappt die Bookmarks in Acrobat aus
colorlinks=true,%
linkcolor=black,%
citecolor=black,%
urlcolor=htworange,%
]{hyperref}

% Nach Dokumentenbeginn ausführen:
\makeatletter
\AtBeginDocument{	
	% pdf Informationen
	\hypersetup{
		pdfauthor={\@author},
		pdftitle={\@title},
		pdfproducer={LaTeX Layout erstellt von Falk-Jonatan Strube}
	}
	% signierte Fußzeile
	%\providecommand{\ofootname}{Mitschrift von}
	\providecommand{\pathtomaster}{../}
	\ofoot{%
	%\ifdefined\ofoottext
	%\else
	%\begin{minipage}{40mm}
	%\begin{flushright}
	%\textcolor{gray}{
	%\ofootname\\
	%Falk-Jonatan Strube\vspace*{-1mm}}
	%\end{flushright}
	%\end{minipage}
	%\fi%
	%\hspace*{1mm}%
	\begin{minipage}{10 mm}
	\includegraphics[height=9mm]{\pathtomaster/FJS-logo_gray}
	\end{minipage}%
	}
}
\makeatother



%%% ==================================================
%%% ===================== MACROS =====================
%%% ==================================================
% Horizontale Linie:
\newcommand{\HRule}[1][\medskipamount]{\par
  \vspace*{\dimexpr-\parskip-\baselineskip+#1}
  \noindent\rule[0.2ex]{\linewidth}{0.2mm}\par
  \vspace*{\dimexpr-\parskip-.5\baselineskip+#1}}
% Gestrichelte horizontale Linie:
\RequirePackage{dashrule}
\newcommand{\HDRule}[1][\medskipamount]{\par
  \vspace*{\dimexpr-\parskip-\baselineskip+#1}
  \noindent\hdashrule[0.2ex]{\linewidth}{0.2mm}{1mm} \par
  \vspace*{\dimexpr-\parskip-.5\baselineskip+#1}}
% Mathe in Anführungszeichen:
\newsavebox{\mathbox}\newsavebox{\mathquote}
\makeatletter
\newcommand{\mq}[1]{% \mathquotes{<stuff>}
  \savebox{\mathquote}{\text{"}}% Save quotes
  \savebox{\mathbox}{$\displaystyle #1$}% Save <stuff>
  \raisebox{\dimexpr\ht\mathbox-\ht\mathquote\relax}{"}#1\raisebox{\dimexpr\ht\mathbox-\ht\mathquote\relax}{''}
}
\makeatother

% Paragraph mit Zähler (Section-Weise)
\newcounter{cparagraphC}
\newcommand{\cparagraph}[1]{
\stepcounter{cparagraphC}
\paragraph{\thesection{}-\thecparagraphC{} #1}
%\addcontentsline{toc}{subsubsection}{\thesection{}-\thecparagraphC{} #1}
\label{\thesection-\thecparagraphC}
}
\makeatletter
\@addtoreset{cparagraphC}{section}
\makeatother


% (Vorlesungs-)Folien einbinden:
\makeatletter
% Folien-Pfad (inkl. Dateiname ohne Endung und ggf. ohne Nummerierung)
\newcommand{\@slidePath}{}
\newcommand{\slidePath}[1]{\renewcommand{\@slidePath}{#1}}
% Folien Eigenschaften
\newcommand{\@slideScale}{0.5}
\newcommand{\slideScale}[1]{\renewcommand{\@slideScale}{#1}}
\AtBeginDocument{
% Folien von einer Datei skaliert
\newcommand{\slide}[2][\@slideScale]{\slides[#1]{}{#2}}
\newcommand{\slideTrim}[6][\@slideScale]{\slides[#1 , clip,  trim = #5cm #4cm #6cm #3cm]{}{#2}}
% Folien von mehreren nummerierten Dateien skaliert
\newcommand{\slides}[3][\@slideScale]{\begin{center}
\includegraphics[page=#3, scale=#1]{\@slidePath #2.pdf}
\end{center}}
}
\makeatother

% \emph{} anders definieren
\makeatletter
\DeclareRobustCommand{\em}{%
  \@nomath\em \if b\expandafter\@car\f@series\@nil
  \normalfont \else \scshape \fi}
\makeatother

% unwichtiges
\newcommand{\unimptnt}[1]{{\transparent{0.5}#1}}

% alph. enumerate
\newenvironment{anumerate}{\begin{enumerate}[label=(\alph*)]}{\end{enumerate}} % Alphabetische Aufzählung

% Hanging parameters
\newcommand{\hangpara}[1]{\par\noindent\hangindent+2em\hangafter=1 #1\par\noindent}

%% EINFACHE BEFEHLE

% Abkürzungen Mathe
\newcommand{\EE}{\mathbb{E}}
\newcommand{\QQ}{\mathbb{Q}}
\newcommand{\RR}{\mathbb{R}}
\newcommand{\CC}{\mathbb{C}}
\newcommand{\NN}{\mathbb{N}}
\newcommand{\ZZ}{\mathbb{Z}}
\newcommand{\PP}{\mathbb{P}}
\renewcommand{\SS}{\mathbb{S}}
\newcommand{\cA}{\mathcal{A}}
\newcommand{\cB}{\mathcal{B}}
\newcommand{\cC}{\mathcal{C}}
\newcommand{\cD}{\mathcal{D}}
\newcommand{\cE}{\mathcal{E}}
\newcommand{\cF}{\mathcal{F}}
\newcommand{\cG}{\mathcal{G}}
\newcommand{\cH}{\mathcal{H}}
\newcommand{\cI}{\mathcal{I}}
\newcommand{\cJ}{\mathcal{J}}
\newcommand{\cM}{\mathcal{M}}
\newcommand{\cN}{\mathcal{N}}
\newcommand{\cP}{\mathcal{P}}
\newcommand{\cR}{\mathcal{R}}
\newcommand{\cS}{\mathcal{S}}
\newcommand{\cZ}{\mathcal{Z}}
\newcommand{\cL}{\mathcal{L}}
\newcommand{\cT}{\mathcal{T}}
\newcommand{\cU}{\mathcal{U}}
\newcommand{\cX}{\mathcal{X}}
\newcommand{\cV}{\mathcal{V}}
\renewcommand{\phi}{\varphi}
\renewcommand{\epsilon}{\varepsilon}
\renewcommand{\theta}{\vartheta}

% Verschiedene als Mathe-Operatoren
\DeclareMathOperator{\arccot}{arccot}
\DeclareMathOperator{\arccosh}{arccosh}
\DeclareMathOperator{\arcsinh}{arcsinh}
\DeclareMathOperator{\arctanh}{arctanh}
\DeclareMathOperator{\arccoth}{arccoth} 
\DeclareMathOperator{\var}{Var} % Varianz 
\DeclareMathOperator{\cov}{Cov} % Co-Varianz 

% Farbdefinitionen
\definecolor{red}{RGB}{180,0,0}
\definecolor{green}{RGB}{75,160,0}
\definecolor{blue}{RGB}{0,75,200}
\definecolor{orange}{RGB}{255,128,0}
\definecolor{yellow}{RGB}{255,245,0}
\definecolor{purple}{RGB}{75,0,160}
\definecolor{cyan}{RGB}{0,160,160}
\definecolor{brown}{RGB}{120,60,10}

\definecolor{itteny}{RGB}{244,229,0}
\definecolor{ittenyo}{RGB}{253,198,11}
\definecolor{itteno}{RGB}{241,142,28}
\definecolor{ittenor}{RGB}{234,98,31}
\definecolor{ittenr}{RGB}{227,35,34}
\definecolor{ittenrp}{RGB}{196,3,125}
\definecolor{ittenp}{RGB}{109,57,139}
\definecolor{ittenpb}{RGB}{68,78,153}
\definecolor{ittenb}{RGB}{42,113,176}
\definecolor{ittenbg}{RGB}{6,150,187}
\definecolor{itteng}{RGB}{0,142,91}
\definecolor{ittengy}{RGB}{140,187,38}

% Textfarbe ändern
\newcommand{\tred}[1]{\textcolor{red}{#1}}
\newcommand{\tgreen}[1]{\textcolor{green}{#1}}
\newcommand{\tblue}[1]{\textcolor{blue}{#1}}
\newcommand{\torange}[1]{\textcolor{orange}{#1}}
\newcommand{\tyellow}[1]{\textcolor{yellow}{#1}}
\newcommand{\tpurple}[1]{\textcolor{purple}{#1}}
\newcommand{\tcyan}[1]{\textcolor{cyan}{#1}}
\newcommand{\tbrown}[1]{\textcolor{brown}{#1}}

% Umstellen der Tabellen Definition
\newcommand{\mpb}[1][.3]{\begin{minipage}{#1\textwidth}\vspace*{3pt}}
\newcommand{\mpe}{\vspace*{3pt}\end{minipage}}

\newcommand{\resultul}[1]{\underline{\underline{#1}}}
\newcommand{\parskp}{$ $\\}	% new line after paragraph
\newcommand{\parnl}{\unskip\parfillskip 0pt \par} % justified new line (brace line!)

% mathematische Abkürzungen
\newcommand{\corr}{\;\widehat{=}\;}
\newcommand{\mdeg}{^{\circ}}
\newcommand{\nok}[2]{\binom{#1}{#2}}	% n über k BESSER: \binom{n}{k}
\newcommand{\mtr}[1]{\begin{pmatrix}#1\end{pmatrix}}	% Matrix
\newcommand{\dtr}[1]{\begin{vmatrix}#1\end{vmatrix}}	% Determinante (Betragsmatrix)
\LetLtxMacro{\originalVec}{\vec}
\renewcommand{\vec}[1]{\underline{#1}}	% Vektorschreibweise
\newcommand{\imptnt}[1]{\colorbox{red!30}{#1}}	% Wichtiges
\newcommand{\intd}[1]{\,\mathrm{d}#1}
\newcommand{\diffd}[1]{\mathrm{d}#1}
% für Module-Rechnung: \pmod{}
\newcommand{\unit}[1]{\,\mathrm{#1}}
%\LetLtxMacro{\ntilde}{\tilde}
%\renewcommand{\tilde}{\widetilde}

% weiter Abkürzungen
\newcommand{\gdw}{genau dann wenn}
\newcommand{\lecdate}[1]{\marginpar{\small\textcolor{gray}{Vorlesung\newline #1}}}

\faculty{Fakultät Informatik/Mathematik}
\chair{Lehrstuhl für Rechnernetze/Kommunikationssysteme}
\subject{lecture}
\title{Internet-Technologien 2}
\author{Falk-Jonatan Strube}
\professor{Prof. Dr.-Ing. Jörg Vogt}

\slidePath{Vorlesung/}
\slideScale{.85}
\makeatletter
\renewcommand{\@pnumwidth}{1.75em} 
\makeatother
\RedeclareSectionCommand[tocnumwidth=4.5em]{subsubsection}
\RedeclareSectionCommand[tocnumwidth=3.5em]{subsection}
\RedeclareSectionCommand[tocnumwidth=2.7em]{section}
\begin{document}
\maketitle
\tableofcontents

\chapter*{Vorbemerkung}
% Nutzer: it2		PW: VLIT2

\chapter{Grundlagen Echtzeitdienste}
\lecdate{17.10.2017}
\section{Motivation}
\slides{it2-10-echtzeitdienste_print}{3}
\subsection{Beispiel IPTV}
\slides{it2-10-echtzeitdienste_print}{4}
\subsubsection*{Probleme IPTV}
\slides{it2-10-echtzeitdienste_print}{5}
\subsection{Echtzeitsysteme}
\slides{it2-10-echtzeitdienste_print}{6}

\section{Verzögerungen in paketorientierten Netzen}
\slides{it2-10-echtzeitdienste_print}{7}
\slides{it2-10-echtzeitdienste_print}{8}
\subsubsection*{Beispiele}
\slides{it2-10-echtzeitdienste_print}{9}
\subsection{Datenrate-Verzögerungsprodukt}
\slides{it2-10-echtzeitdienste_print}{10}
\subsection{Round-Trip-Time (RTT)}
\slides{it2-10-echtzeitdienste_print}{11}
\subsubsection*{RTT-Bestimmung}
\slides{it2-10-echtzeitdienste_print}{13}
\subsubsection{Retransmission-Timer bei TCP}
\slides{it2-10-echtzeitdienste_print}{12}
\subsubsection*{Bestimmung des RTO bei TCP}
\slides{it2-10-echtzeitdienste_print}{14}
\subsubsection{Latenz von TCP: Fast Retransmit}
\slides{it2-10-echtzeitdienste_print}{15}
\subsection{Anforderungen für VoIP}
\slides{it2-10-echtzeitdienste_print}{16}
\subsubsection*{Verzögerungen}
\slides{it2-10-echtzeitdienste_print}{17}
Häufigster Grund für Verzögerung/Paketverlust: Pufferüberlauf (am Router)!
\subsubsection*{Jitter-Puffer}
\slides{it2-10-echtzeitdienste_print}{18}
\subsubsection*{Protokolle im Umfeld VoIP}
\slides{it2-10-echtzeitdienste_print}{19}

\section{Quality of Service}
\slides{it2-10-echtzeitdienste_print}{22}
\subsection{QoS Schicht 3}
\slides{it2-10-echtzeitdienste_print}{20}
\subsection{QoS Schicht 2}
\slides{it2-10-echtzeitdienste_print}{21}

\section{Fehlerschutz}
\lecdate{24.10.2017}
\slides{it2-10-echtzeitdienste_print}{23}
\subsection{ARQ}
\slides{it2-10-echtzeitdienste_print}{24}
\subsection{FEC}
\slides{it2-10-echtzeitdienste_print}{25}
\subsection{3}
\slides{it2-10-echtzeitdienste_print}{26}

\section{Problemfelder}
\subsection{Unicast vs Multicast}
\slides{it2-10-echtzeitdienste_print}{27}
\subsection{Datenkarussell}
\slides{it2-10-echtzeitdienste_print}{28}
\subsection{Datenübertragung in P2P-Netzen}
\slides{it2-10-echtzeitdienste_print}{29}
\subsection{Mobile Empfänger}
\slides{it2-10-echtzeitdienste_print}{30}
\subsection{IP-Multicast über WLAN}
\slides{it2-10-echtzeitdienste_print}{31}
Als Empfänger hat man nur macht über Schicht 7: Dem entsprechend kann der Fehlerschutz nur dort über FEC passieren.
\subsection{Application-Layer Hybrid-Fehlerschutzverfahren}
\slides{it2-10-echtzeitdienste_print}{32}

\section{Zusammenfassung}
\slides{it2-10-echtzeitdienste_print}{34}
%\subsubsection*{Ziele der Vorlesung}
%\slides{it2-10-echtzeitdienste_print}{35}




\chapter{Streaming}
\section{Unterteilung der Streamingverfahren}
\subsection{Unterteilung nach Anwendnug}
\slides{it2-15-streaming_print}{3}
\subsection{Unterteilung nach Techniken}
\slides{it2-15-streaming_print}{4}
\subsection{Streaming mittels Webserver}
\slides{it2-15-streaming_print}{5}
\subsection{Dedizierte Streaming Protokolle}
\slides{it2-15-streaming_print}{6}
\subsection{Progressiver Download}
\slides{it2-15-streaming_print}{7}
\subsection{HTML mit Video-Tag}
\slides{it2-15-streaming_print}{8}
%%%% Pause

\chapter{RTSP}
\lecdate{24.10.2017-2}
\section{Grundlagen}
\slides{it2-17-rtsp_print}{3}
\subsection{Out-of-Band-Steuerung}
\slides{it2-17-rtsp_print}{4}
\subsection{Ports}
\slides{it2-17-rtsp_print}{5}
\subsection{Beispielszenario}
\slides{it2-17-rtsp_print}{6}
\subsection{Beispiel einer Metadatei}
\slides{it2-17-rtsp_print}{7}
\subsection{Methoden}
\slides{it2-17-rtsp_print}{8}
\subsection{Ablauf}
\slides{it2-17-rtsp_print}{9}

\section{Zustände}
\slides{it2-17-rtsp_print}{10}
\subsection{Statuszeile und Codes}
\slides{it2-17-rtsp_print}{11}
\subsection{Beispielablauf}
\slides{it2-17-rtsp_print}{12}

\section{Minimaler Server}
\slides{it2-17-rtsp_print}{13}
\subsection{Anforderungen}
\slides{it2-17-rtsp_print}{14}

\section{Methoden}
\subsection{SETUP}
\slides{it2-17-rtsp_print}{15}
\subsection{PLAY}
\slides{it2-17-rtsp_print}{16}
\subsection{PAUSE}
\slides{it2-17-rtsp_print}{17}
\subsection{DESCRIBE}
\slides{it2-17-rtsp_print}{18}

\section{Session Description Protocol}
\slides{it2-17-rtsp_print}{19}
\subsubsection*{Beispiel}
\slides{it2-17-rtsp_print}{20}

\section{Zusammenfassung}
\slides{it2-17-rtsp_print}{23}













\chapter{RTP}
\lecdate{07.11.2017-2}
\section{Grundlagen}
\slides{it2-18-rtp_print}{3}
\subsection{Architektur}
\slides{it2-18-rtp_print}{4}
\subsubsection*{RTP-Felder}
\slides{it2-18-rtp_print}{5}
\subsection{Payload (Beispiel)}
\slides{it2-18-rtp_print}{6}
\subsection{Header-Aufbau}
\slides{it2-18-rtp_print}{7}
\section{Programmierung in Java}
\subsection{Bitoperationen}
\slides{it2-18-rtp_print}{8}
\begin{align*}
20 << 2 &= 80\\
20 >> 2 &= 5\\
-1 >> 2 &= \\
-1 >>> 2 &=
\end{align*}
\subsection{Setzen des Headers}
\slides{it2-18-rtp_print}{9}

\section{RealTime Control Protocol (RTCP)}
\slides{it2-18-rtp_print}{10}

\section{Zusammenfassung}
\slides{it2-18-rtp_print}{11}








\chapter{FEC}
\section{Kanalmodelle}
\subsection{BSC und BEC}
\slides{it2-20-fec-einfuehrung_print}{3}
Anmerkung: Durch Prüfsummen kommen auf der Anwendungsebene in der Regel keine Pakete mit Bitfehler an (werden vom WLAN-Chip o.ä. abgefangen). Deswegen BSC hier weniger relevant.\\
Es kann aber passieren, dass Bits verloren gehen (weniger Bitweise, sondern mehr Paketweise) $\To$ BEC.
\subsection{Additive white Gaussian Noise (AWGN)}
\slides{it2-20-fec-einfuehrung_print}{4}
Kann bspw. beim Auslesen einer (klassischen) Festplatte vorkommen.
\slides{it2-20-fec-einfuehrung_print}{5}






\chapter{Blockcodierung - Begriffe}
\section{Codierungsverfahren}
\slides{it2-30-blockcodes_print}{3}
\section{Begriffe}
\slides{it2-30-blockcodes_print}{4}
\subsection{Systematische und Nichtsystematische Codierung}
\slides{it2-30-blockcodes_print}{5}
\subsection{Hamming-Distanz und Hamming-Gewicht}
\slides{it2-30-blockcodes_print}{6}
\slides{it2-30-blockcodes_print}{7}
\subsection{Fehlererkennung}
\slides{it2-30-blockcodes_print}{8}
\subsection{Fehlerkorrektur}
\slides{it2-30-blockcodes_print}{9}
\subsection{Fehlererkennung und -korrektur}
\slides{it2-30-blockcodes_print}{10}
\subsection{Auslöschungskorrektur}
\slides{it2-30-blockcodes_print}{11}
\section{Notation von Blocks}
\slides{it2-30-blockcodes_print}{12}
\section{Übersicht der FEC-Verfahren}
\slides{it2-30-blockcodes_print}{13}
\subsection{Restfehlerwahrscheinlichkeit bei Fehlerkorrektur}
\lecdate{28.11.2017}
\slides{it2-30-blockcodes_print}{14}
\slides{it2-30-blockcodes_print}{15}
\section{Codiergewinn}
\subsection{Vergleich codierter mit uncodierter Übertragung}
\slides{it2-30-blockcodes_print}{16}
\unimptnt{
\subsection{Codiergewinn bei AWGN-Kanal}
\slides{it2-30-blockcodes_print}{17}
\subsection{Codiergewinn am Beispiel des Hamming-Codes}
\slides{it2-30-blockcodes_print}{18}
}
\subsection{Vergleich von Auslöschungscodes}
\slides{it2-30-blockcodes_print}{19}
\subsection{Decodierergebnisse}
\slides{it2-30-blockcodes_print}{20}
\section{Decodierprinzipien}
\slides{it2-30-blockcodes_print}{21}
\subsection{Bounded Minimum Distance (BMD)}
\slides{it2-30-blockcodes_print}{22}
\subsection{Maximum Likelihood (MLD)}
\slides{it2-30-blockcodes_print}{23}
Es ist nicht mit abgedeckt, wie oft (und damit wie Wahrscheinlich) ein Codewort ist, nur die Nähe wird berücksichtigt. Dafür benutzt man MAP:
\subsection{Maximum Aposteriori (MAP)}
\slides{it2-30-blockcodes_print}{24}
Wird in der Regel nicht benötigt, da bei Video-codierung usw. verfahren angewendet werden, dass die zu übertragenden Pakete immer ein ähnliches statistisches Aufkommen haben.

\subsection{Decodierverfahren (Soft-/Hard-Decision)}
\slides{it2-30-blockcodes_print}{25}









\chapter{Blockcodierung - Lineare Codes}
\section{Motivation}
\subsection{Dilemma der Kanalcodierung}
\slides{it2-35-lincodes_print}{3}
\subsection{Motivation für Lineare Codes}
\slides{it2-35-lincodes_print}{4}
\subsection{Vektorraum}
\slides{it2-35-lincodes_print}{5}
GF: Galois-Feld (endlicher Körper)… (endlicher) Körper: man kann jede Operation ausführen (Multiplikation/Addition).
\subsubsection*{Beispiel Vektorraum mit GF(3)}
\slides{it2-35-lincodes_print}{6}

\section{Definition Linearer Code}
\slides{it2-35-lincodes_print}{7}
(2. Beispiel: man kommt nie auf $111$ $\to$ ist also kein linearer Code)\\
Hinweis: auf $000$ kommt man immer, wenn man einen beliebiges Codewort mit sich selbst addiert.
\subsection{Hamming-Gewicht}
\slides{it2-35-lincodes_print}{8}
Der Abstand zwischen zwei Codewörtern ist wieder ein Codewort. Also muss man nur die einsen Zählen…
\subsection{Gewichtsverteilung}
\slides{it2-35-lincodes_print}{9}
\subsection{Hamming-Codes}
\slides{it2-35-lincodes_print}{10}
Rote Codewörter: „Einheitsvektoren“ (haben jeweils nur eine $1$) $\to$ durch Skalierung, … kann jeder anderer „Vektor“/Codewort dargestellt werden.
\subsection{Singleton-Schranke}
\slides{it2-35-lincodes_print}{11}

\section{Hamming-Schranke}
\slides{it2-35-lincodes_print}{12}
\subsubsection*{Beispiel}
\slides{it2-35-lincodes_print}{13}

\section{Generatormatrix}
\slides{it2-35-lincodes_print}{14}
\subsection{Bildung des Codewortes}
\slides{it2-35-lincodes_print}{15}
\subsection{Systematische Codierung}
\slides{it2-35-lincodes_print}{16}
\subsection{Codierschaltung}
\slides{it2-35-lincodes_print}{17}

\section{Dualer Code}
\slides{it2-35-lincodes_print}{18}

\section{Prüfmatrix}
\slides{it2-35-lincodes_print}{19}
\subsection*{Kanal}
\slides{it2-35-lincodes_print}{20}

\section{Nebenklassenzerlegung}
\slides{it2-35-lincodes_print}{21}
Nebenklassenführer: Der Code der Nebenklasse mit den wenigsten einsen (in Fällen 0-5 eindeutig, bei 6 und 7 nicht eindeutig).
\subsubsection*{Beispiel}
\slides{it2-35-lincodes_print}{22}
Nebenklassencodes sind willkürlich gewählt (solange sie den Voraussetzungen entsprechen).
\subsection{Fehlerkorrektur}
\slides{it2-35-lincodes_print}{23}

\section{Syndromdekodierung}
\lecdate{05.12.2017}
\subsection{Fehlererkennung}
\slides{it2-35-lincodes_print}{24}
\subsubsection*{Beispiele}
\slides{it2-35-lincodes_print}{25}
\subsection{Fehlerkorrektur}
\slides{it2-35-lincodes_print}{26}
Empfänger erhält $y=x+e$
\begin{enumerate}
\item $s=y\cdot H^T$ $\to$ $\hat{e}$
\item $\hat{x}=y+\hat{e}$
\end{enumerate}
\subsubsection*{Syndromtabelle}
\slides{it2-35-lincodes_print}{27}

\section{Hamming Codes}
\subsection{Definition}
\slides{it2-35-lincodes_print}{28}
\subsection{Erzeugung}
\slides{it2-35-lincodes_print}{29}
\subsection{Decodierung}
\slides{it2-35-lincodes_print}{30}
\subsubsection*{Eigenschaften Syndromdecodierung}
\slides{it2-35-lincodes_print}{31}
\subsection{Beispiele}
\subsubsection*{Hamming-Code (63,57,3)}
\slides{it2-35-lincodes_print}{32}
\subsubsection*{BCH-Code (511,259,61)}
\slides{it2-35-lincodes_print}{33}

\section{Zusammenfassung}
\slides{it2-35-lincodes_print}{34}




\chapter{Blockcodierung - Zyklische Codes}
\section{Einführung}
\slides{it2-40-zykcodes_print}{3}
\subsection{Definition zyklischer Code}
\slides{it2-40-zykcodes_print}{4}
\subsection{Beschreibung zyklischer Code}
\slides{it2-40-zykcodes_print}{5}
\subsection{Polynomdivision}
\slides{it2-40-zykcodes_print}{6}

\section{Zyklische Verschiebung}
\slides{it2-40-zykcodes_print}{7}
\subsubsection*{Beispiel}
\slides{it2-40-zykcodes_print}{8}

\section{Generatorproblem}
\slides{it2-40-zykcodes_print}{9}
\slides{it2-40-zykcodes_print}{10}
\subsubsection*{Beispiel}
\slides{it2-40-zykcodes_print}{11}

\section{Algorithmus für systematische Codierung}
\slides{it2-40-zykcodes_print}{12}
\subsubsection*{Beispiel}
\slides{it2-40-zykcodes_print}{13}

\section{Dekodierung}
\slides{it2-40-zykcodes_print}{14}
% Folien Fehlerkorrektur unwichtig

\section{CRC-Codes}
\slides{it2-40-zykcodes_print}{15}












\chapter{Blockcodierung - Galois-Felder}
\lecdate{19.12.2017}
\section{Einführung}
\slides{it2-50-galois_print}{3}
\subsection{Definition Galois-Feld}
\slides{it2-50-galois_print}{4}
\subsection{Definition Prim- und Erweiterungskörper}
\slides{it2-50-galois_print}{5}
\subsubsection*{Beispiel GF(3)}
\slides{it2-50-galois_print}{6}
\subsubsection*{Beispiel Nichtexistenz von GF(4)}
\slides{it2-50-galois_print}{7}
\subsubsection*{Beispiel GF(5)}
\slides{it2-50-galois_print}{8}

\subsection{Primitives Element}
\slides{it2-50-galois_print}{9}
\subsubsection*{Ordnung eines Elements}
\slides{it2-50-galois_print}{10}
\subsection{Irreduzible Polynome}
\slides{it2-50-galois_print}{11}
\subsection{Erweiterungskörper}
\slides{it2-50-galois_print}{12}
\subsubsection{Existenz}
\slides{it2-50-galois_print}{13}
\subsubsection{Konstruktion}
\slides{it2-50-galois_print}{14}
\subsubsection{Rechnung}
\slides{it2-50-galois_print}{15}
\subsection{Primitives Polynom}
\slides{it2-50-galois_print}{16}
\slides{it2-50-galois_print}{17}
\subsection{Alternative Berechnung Modulo}
\slides{it2-50-galois_print}{18}
\subsubsection*{Beispiel GF(4)}
\slides{it2-50-galois_print}{19}
\slides{it2-50-galois_print}{20}

\section{Darstellungsarten der Erweiterungskörper}
\slides{it2-50-galois_print}{21}

\subsection{Beispiel GF(8)}
\slides{it2-50-galois_print}{22}
\slides{it2-50-galois_print}{23}
\slides{it2-50-galois_print}{24}
\slides{it2-50-galois_print}{25}

\subsection{Beispiel GF(16)}
\slides{it2-50-galois_print}{26}

\section{Implementierung}
\slides{it2-50-galois_print}{27}
\subsubsection*{Beispiel}
\slides{it2-50-galois_print}{28}












%\newpage
%\printbibliography

\end{document}